\scanpage{52}
\margin{\unsure{1}}
\begin{center}4 - 38 -\hfill n\textsuperscript{o} 392\end{center}

où $I$ est pseudo-nul. Il résulte donc du n\textsuperscript{o} 7, prop. 16, (v) et cor. 1 de la prop. 16, que l'on a $c_B(M\otimes_A B)=c_B(Q\otimes_A B)=-c_B(\unsure{\chi|\gamma}_B(Q\otimes_A B))=-c_B(i(\unsure{\chi|\gamma}_A(Q)))$ en vertu de (ii) ; mais par définition de l'homomorphisme $i:C(A)\longrightarrow C(B)$, on a $c_B(i(\unsure{\chi|\gamma}_A(Q)))=i(c_A(\unsure{\chi|\gamma}_A(Q)))=-i(c_A(M))$, ce qui achève de prouver (23).

\uline{Remarques}.- 1) Si $M$ est un $A$-module réflexif, $M\otimes_A B$ n'est pas nécessairement réflexif, (exerc.\quad). Il en est toutefois ainsi lorsque $B$ est un $A$-module \uline{plat}, (n\textsuperscript{o} 2, prop. 8).

2) Soit $C$ un troisième anneau noethérien intégralement clos, tel que $B\subset C$ et que $C$ soit un $B$-module de type fini, (donc aussi un $A$-module de type fini). On a alors les formules de transitivité :
\[
\begin{aligned}
(25)\qquad &N_{C/A}=N_{B/A}\circ N_{C/B},\\
(26)\qquad &i_{C/A}=i_{C/B}\circ i_{B/A}
\end{aligned}
\]
qui résultent immédiatement des formules de transitivité pour les indices de ramification et les degrés résiduels, (chap. VI, § 8, n\textsuperscript{o} 1, lemme 1).

9. \uline{Un théorème de réduction}.

Les notations et hypothèses sont de nouveau celles des n\textsuperscript{os} 2 à 7.

LEMME 6.- \uline{Soient} $R$ \uline{un anneau commutatif}, $\underline{p}_i$, $(1\le i\le n)$, \uline{des idéaux premiers de} $R$, \uline{deux à deux distincts}.

(i) \uline{Pour} $1\le i\le n$, \uline{soit} $H_i$ \uline{une partie de} $R/\underline{p}_i$ \uline{vérifiant la condition suivante} : \uline{il n'existe aucun élément} $\alpha_i\in R/\underline{p}_i$ \uline{tel que} $\alpha_i+H_i$ \uline{contienne un idéal} $\ne0$ \uline{de} $R/\underline{p}_i$. \uline{Alors il existe} $a\in R$ \uline{tel que}, \uline{pour} $1\le i\le n$, \uline{l'image canonique de} $a$ \uline{dans} $R/\underline{p}_i$ \uline{n'appartienne pas à} $H_i$.

(ii) \uline{Si} $\operatorname{Card}(H_i)<\operatorname{Card}(R/\underline{p}_i)$, \uline{les} $H_i$ \uline{vérifient la condition de} (i).

(i) Raisonnons par récurrence sur $n$, le cas $n=0$ étant trivial. Soit donc $n\ge1$. Quitte à faire une permutation sur les indices $i$, on peut supposer que $\underline{p}_1$ est

\margin{51}
